\subsection{Why a Zero Dot Product Means Perpendicularity}

The condition
\[
\mathbf u\cdot\mathbf v=0
\]
should not be accepted as an unexplained rule.  It is already hidden in the
coordinate criterion for perpendicular lines established in the preceding
chapter.

Take two nonzero vectors
\[
\mathbf u=(u_1,u_2),
\qquad
\mathbf v=(v_1,v_2)
\]
and choose an arbitrary point
\[
P=(x_0,y_0).
\]
Use the vectors as coordinate displacements from that same point:
\[
Q=P+\mathbf u=(x_0+u_1,\,y_0+u_2),
\]
\[
R=P+\mathbf v=(x_0+v_1,\,y_0+v_2).
\]
The line through $P$ and $Q$ has direction $\mathbf u$, while the line through
$P$ and $R$ has direction $\mathbf v$.

\begin{center}
\begin{tikzpicture}[scale=0.9]
    \coordinate (P) at (0.7,0.6);
    \coordinate (Q) at (4.5,2.1);
    \coordinate (R) at (-0.5,3.6);
    \draw[subset,->] (P) -- (Q) node[midway,below right] {$\mathbf u$};
    \draw[subset,->] (P) -- (R) node[midway,left] {$\mathbf v$};
    \draw[helper] (-0.25,0.225) -- (5.45,2.475);
    \draw[helper] (1.0,-0.15) -- (-0.8,4.35);
    \fill[geometricSubsetColor] (P) circle (2pt) node[below left] {$P$};
    \fill[geometricSubsetColor] (Q) circle (2pt) node[above right] {$Q$};
    \fill[geometricSubsetColor] (R) circle (2pt) node[above left] {$R$};
\end{tikzpicture}
\end{center}

First suppose that neither line is vertical, so
\[
u_1\neq0,
\qquad
v_1\neq0.
\]
Their slopes are obtained directly from the coordinate differences:
\[
m_{\mathbf u}
=
\frac{(y_0+u_2)-y_0}{(x_0+u_1)-x_0}
=
\frac{u_2}{u_1},
\]
\[
m_{\mathbf v}
=
\frac{(y_0+v_2)-y_0}{(x_0+v_1)-x_0}
=
\frac{v_2}{v_1}.
\]
The previously proved criterion for perpendicular nonvertical lines is
\[
m_{\mathbf u}m_{\mathbf v}=-1.
\]
Substitution gives
\[
\frac{u_2}{u_1}\frac{v_2}{v_1}=-1.
\]
Multiplying by $u_1v_1$ yields
\[
u_2v_2=-u_1v_1,
\]
and therefore
\[
u_1v_1+u_2v_2=0.
\]
But the left-hand side is exactly the coordinate definition of the dot product:
\[
\mathbf u\cdot\mathbf v=u_1v_1+u_2v_2.
\]
Hence, in the nonvertical case,
\[
\text{the two direction lines are perpendicular}
\iff
\mathbf u\cdot\mathbf v=0.
\]

The remaining case is a horizontal line perpendicular to a vertical line.  A
nonzero horizontal direction has the form
\[
\mathbf u=(u_1,0),
\]
while a nonzero vertical direction has the form
\[
\mathbf v=(0,v_2).
\]
Then
\[
\mathbf u\cdot\mathbf v
=
u_1\cdot0+0\cdot v_2
=0.
\]
Conversely, if $\mathbf u=(u_1,0)$ is nonzero and
\[
\mathbf u\cdot\mathbf v=u_1v_1=0,
\]
then $v_1=0$, so $\mathbf v$ is vertical.

\begin{theorembox}
For nonzero vectors $\mathbf u,\mathbf v\in\mathbb R^2$, the lines through a
common point with directions $\mathbf u$ and $\mathbf v$ are perpendicular if
and only if
\[
\mathbf u\cdot\mathbf v=0.
\]
\end{theorembox}

The choice of the common point $P$ was arbitrary.  It disappears from every
slope calculation because only coordinate differences remain.  Thus
perpendicularity depends on the two displacement vectors, not on where their
representatives are drawn.

\begin{examplebox}
Take
\[
\mathbf u=(3,2),
\qquad
\mathbf v=(-2,3).
\]
The corresponding slopes are
\[
m_{\mathbf u}=\frac23,
\qquad
m_{\mathbf v}=-\frac32,
\]
so their product is $-1$.  The same fact appears immediately in the dot
product:
\[
\mathbf u\cdot\mathbf v
=3(-2)+2(3)
=0.
\]
The algebraic test and the geometric right-angle condition are therefore the
same statement written in two languages.
\end{examplebox}